\section{Accessibility is a receiving operation}
The domain in the phrase ``accessible record'' must be part of the experiment. A record may be perfectly distinguishable jointly and unreadable from a restricted subsystem. A signal may already have been emitted but not yet have reached a receiver. A memory may survive while its decoder is unavailable. These cases have different operational meanings.

For any two candidate memory states, a trace-preserving physical access channel $\mathcal A$ cannot increase their trace distance:
\begin{equation}
 \tfrac12\norm{\mathcal A(\rho_0)-\mathcal A(\rho_1)}_1
 \leq\tfrac12\norm{\rho_0-\rho_1}_1.
\end{equation}
The proof is the variational characterization of trace distance: every allowed output POVM pulls back through the unital dual channel to an allowed input POVM. An inaccessible fragment can reduce readout performance; calling it inaccessible does not remove its contribution to the global state or restore a lost marginal fringe.

Chapter V adds a causal constraint to this account. An exterior record must reach an exterior receiver through an allowed physical process. Late receipt of signals from a finite source-age interval does not give the emitter an infinite history. A boundary convention that excludes a completed parent-ledger state supplies neither a quantum channel across that boundary nor a rule selecting one outcome. The state, field dynamics and receiving map remain necessary.

The finite tensor-product examples in this chapter are operational models of separated systems. They do not establish an exact factorization of a gravitational or constrained gauge-field Hilbert space across a horizon. If a proposed theory replaces that subsystem structure by an observable algebra, its accessible operations and probabilities must be specified there instead.

These restrictions preserve the source's physical distinction: a write is more than an unresolved possibility. They also prevent the word ``write'' from performing work that belongs to a Hamiltonian, an instrument, a decoder or a causal signal.


